\documentclass[../../../include/open-logic-section]{subfiles}

\begin{document}

\olfileid[fr]{sth}{z}{sep}
\olsection{Séparation}

Commençons par un principe destiné à remplacer la compréhension naïve :

\begin{axiom}[Schéma de séparation] Pour toute formule $\phi(x)$,
l'énoncé suivant est un axiome : quel que soit $A$, l'ensemble
$\Setabs{x \in A}{\phi(x)}$ existe.
\end{axiom}

Remarquons qu'il ne s'agit pas d'un axiome unique, mais d'un
\emph{schéma} d'axiomes. Il y a une \emph{infinité} d'axiomes de
séparation, un pour chaque formule $\phi(x)$. On peut également écrire
le schéma sous la forme suivante, qui est la forme habituelle :

\begin{defish}
Pour toute formule $\phi(x)$ qui ne contient pas «~$S$~», l'énoncé
suivant est un axiome :
\[
	\forall A \exists S \forall x(x \in S \liff (\phi(x) \land x \in A)).
\]
\end{defish}

Conformément à la convention énoncée au début de la
\olref[sth][][]{part}, les formules~$\phi$ des axiomes de séparation
peuvent contenir des paramètres.\footnote{Pour comprendre ce que cela
signifie, voir le passage immédiatement après la
\olref[sfr][infinite][induction]{natinductionschema}.}

La conception cumulative et itérative des ensembles que nous avons
présentée justifie directement la séparation. Pour le voir, soit $A$
un ensemble. L'ensemble $A$ est donc formé à une certaine étape~$S$
(par \stageshier). Puisque $A$ est formé à l'étape~$S$, tous les
!!{element}s de $A$ ont été formés avant l'étape~$S$ (par \stagesacc).
Considérons en particulier tous les ensembles qui appartiennent à~$A$
et qui satisfont aussi $\phi$ : eux aussi ont manifestement été formés
avant l'étape~$S$. Ils sont donc réunis en un ensemble
$\Setabs{x \in A}{\phi(x)}$ à cette même étape~$S$ (par \stagesacc).

À la différence de la compréhension naïve, ce principe évite le
paradoxe de Russell. Nous ne pouvons pas simplement affirmer
l'existence de l'ensemble $\Setabs{x}{x \notin x}$. En revanche,
\emph{étant donné} un ensemble~$A$, nous pouvons affirmer l'existence
de l'ensemble $R_A = \Setabs{x \in A}{x \notin x}$. Mais tout ce que
cela permet de démontrer, c'est que $R_A \notin R_A$ et
$R_A \notin A$, ce qui n'a rien d'inquiétant.

La séparation a toutefois une conséquence immédiate et frappante :

\begin{thm}\ollabel{thm:NoUniversalSet}
Il n'existe pas d'ensemble \emph{universel}, c'est-à-dire que
$\Setabs{x}{x = x}$ n'existe pas.
\end{thm}

\begin{proof}
Supposons, pour obtenir une contradiction, que $V$ soit un ensemble
universel. Alors, par séparation, l'ensemble
$R = \Setabs{x \in V}{x \notin x} = \Setabs{x}{x \notin x}$ existe.
Comme $V$ contient tous les ensembles, la définition de $R$ donne
$R \in R \liff R \notin R$, ce qui est impossible.
\end{proof}

L'absence d'un ensemble universel, et même le caractère ouvert de la
hiérarchie des ensembles, compte parmi les idées les plus fondamentales
de la conception cumulative et itérative. Il vaut donc la peine de
voir que, intuitivement, on peut y parvenir par une autre voie. Un
ensemble universel doit être !!a{element} de lui-même. Or, selon
notre conception cumulative et itérative, tout ensemble apparaît pour
la première fois dans la hiérarchie à la toute première étape après
l'apparition de tous ses !!{element}s. Il en
résulte qu'\emph{aucun} ensemble ne s'appartient à lui-même : un tel
ensemble devrait apparaître pour la première fois immédiatement après
l'étape où il apparaît pour la première fois, ce qui est absurde.
(Nous verrons à la \olref[sth][spine][rank]{defnsetrank} comment
rendre cette explication plus rigoureuse grâce à la notion de
«~rang~» d'un ensemble. Nous aurons toutefois besoin de quelques
axiomes supplémentaires.)

Voici quelques autres conséquences de la séparation et de
l'extensionnalité.

\begin{prop}\ollabel{prop:emptyexists}
Si au moins un ensemble existe, alors $\emptyset$ existe.
\end{prop}

\begin{proof}
Si $A$ est un ensemble, alors
$\emptyset = \Setabs{x \in A}{x \neq x}$ existe par séparation.
\end{proof}

\begin{prop}
$A \setminus B$ existe quels que soient les ensembles $A$ et $B$.
\end{prop}

\begin{proof}
$A \setminus B = \Setabs{x \in A}{x \notin B}$ existe par séparation.
\end{proof}

Il s'avère aussi que les intersections (presque) arbitraires existent :

\begin{prop}\ollabel{prop:intersectionsexist}
Si $A \neq \emptyset$, alors
$\bigcap A = \Setabs{x}{(\forall y \in A)x \in y}$ existe.
\end{prop}

\begin{proof}
Soit $A \neq \emptyset$ ; il existe donc un $c \in A$. Alors
$\bigcap A = \Setabs{x}{(\forall y \in A)x \in y}
= \Setabs{x \in c}{(\forall y \in A)x \in y}$, qui existe par
séparation.
\end{proof}

Notons toutefois la condition $A \neq \emptyset$ : si elle n'était pas
remplie, $\bigcap \emptyset$ serait l'ensemble universel, par vacuité, ce qui
contredirait le \olref{thm:NoUniversalSet}.

\end{document}
